\documentclass[10pt,a4paper]{article}

\usepackage[utf8]{inputenc}
\usepackage[T1]{fontenc}
\usepackage[french]{babel}
\usepackage[margin=1.5cm]{geometry}
\usepackage{amsmath,amssymb}
\usepackage{xcolor}
\usepackage{listings}
\usepackage{tcolorbox}
\tcbuselibrary{breakable}
\usepackage{enumitem}
\usepackage{array}
\usepackage{booktabs}

% Configuration listings
\lstset{
    language=C,
    basicstyle=\ttfamily\scriptsize,
    backgroundcolor=\color{gray!10},
    keywordstyle=\color{blue},
    commentstyle=\color{gray},
    stringstyle=\color{red},
    numbers=left,
    numberstyle=\tiny,
    breaklines=true,
    frame=single,
    tabsize=4,
    showstringspaces=false,
    inputencoding=utf8,
    extendedchars=true,
    literate={é}{{\'e}}1 {è}{{\`e}}1 {à}{{\`a}}1 {ç}{{\c{c}}}1 
             {ê}{{\^e}}1 {î}{{\^i}}1 {ô}{{\^o}}1 {û}{{\^u}}1 {ù}{{\`u}}1
             {É}{{\'E}}1 {È}{{\`E}}1 {À}{{\`A}}1 {Ç}{{\c{C}}}1
}

% Boxes pédagogiques
\newtcolorbox{exercicebox}[1]{
    colback=green!5,colframe=green!50!black,
    fonttitle=\bfseries,title=#1,breakable
}

\newtcolorbox{infobox}[1]{
    colback=cyan!5,colframe=cyan!75!black,
    fonttitle=\bfseries,title=#1,breakable
}

\title{\textbf{TD Chapitre 2 : Organisation Mémoire}\\
\large Pointeurs, Tableaux et Forensics Mémoire}
\author{Dr. El Hadji Bassirou TOURÉ\\
Département de Mathématiques et Informatique\\
Université Cheikh Anta Diop de Dakar}
\date{}

\begin{document}
\maketitle

%==============================================================================
\begin{exercicebox}{Exercice 1 : Endianness et Représentation Mémoire}
%==============================================================================

\textbf{Partie A :} La valeur \texttt{0x12345678} est stockée à l'adresse \texttt{0x100} sur 4 octets.

\begin{enumerate}[label=\alph*)]
    \item Dessiner le contenu de la mémoire (adresses 0x100 à 0x103) en \textbf{big-endian}.
    \item Dessiner le contenu de la mémoire en \textbf{little-endian}.
    \item Sur une machine little-endian, si on lit l'octet à l'adresse \texttt{0x100}, quelle valeur obtient-on ?
    \item Sur une machine little-endian, si on lit un \texttt{short} à l'adresse \texttt{0x100}, quelle valeur obtient-on ?
\end{enumerate}

\textbf{Partie B :} Voici un dump mémoire d'une machine little-endian :

\begin{center}
\renewcommand{\arraystretch}{1.2}
\begin{tabular}{|c|c|c|c|c|}
\hline
\textbf{Addr} & \textbf{+0} & \textbf{+1} & \textbf{+2} & \textbf{+3} \\
\hline
0x200 & 78 & 56 & 34 & 12 \\
\hline
\end{tabular}
\end{center}

\begin{enumerate}[label=\alph*)]
    \setcounter{enumi}{4}
    \item Quelle est la valeur de l'entier 32 bits stocké à 0x200 ?
    \item Si ce même dump provenait d'une machine big-endian, quelle serait la valeur ?
\end{enumerate}

\end{exercicebox}

%==============================================================================
\begin{exercicebox}{Exercice 2 : Pointeurs et Arithmétique}
%==============================================================================

Soit le code suivant sur une machine 32 bits :

\begin{lstlisting}[numbers=none]
int arr[5] = {10, 20, 30, 40, 50};
/* &arr[0] = 0x1000 */

int *p = arr;
int *q = &arr[3];
\end{lstlisting}

\begin{enumerate}[label=\alph*)]
    \item Quelle est la valeur de \texttt{p} ?
    \item Quelle est la valeur de \texttt{q} ?
    \item Quelle est la valeur de \texttt{p + 2} ?
    \item Quelle est la valeur de \texttt{q - p} ? (attention au type !)
    \item Quelle est la valeur de \texttt{*(p + 4)} ?
    \item Quelle est la valeur de \texttt{p[2]} ? Équivalent à quelle expression ?
    \item Quelle est la valeur de \texttt{(char*)q - (char*)p} ? (résultat en octets)
\end{enumerate}

\textbf{Partie B :} Soit \texttt{void *base = (void *)0x1000;}

Calculer les adresses résultantes :

\begin{center}
\renewcommand{\arraystretch}{1.3}
\begin{tabular}{|l|c|l|}
\hline
\textbf{Expression} & \textbf{sizeof(*p)} & \textbf{Adresse résultante} \\
\hline
\texttt{(int *)base + 3} & 4 & \\
\hline
\texttt{(short *)base + 3} & 2 & \\
\hline
\texttt{(char *)base + 3} & 1 & \\
\hline
\texttt{(double *)base + 2} & 8 & \\
\hline
\end{tabular}
\end{center}

\end{exercicebox}

%==============================================================================
\begin{exercicebox}{Exercice 3 : Chaînes de Caractères}
%==============================================================================

\textbf{Partie A :} Soit le dump mémoire suivant :

\begin{center}
\renewcommand{\arraystretch}{1.2}
\begin{tabular}{|c|c|c|c|c|c|c|c|c|}
\hline
\textbf{Addr} & \textbf{+0} & \textbf{+1} & \textbf{+2} & \textbf{+3} & \textbf{+4} & \textbf{+5} & \textbf{+6} & \textbf{+7} \\
\hline
0x20 & 42 & 6F & 6E & 6A & 6F & 75 & 72 & 00 \\
\hline
0x28 & 55 & 43 & 41 & 44 & 00 & 00 & 00 & 00 \\
\hline
\end{tabular}
\end{center}

\begin{enumerate}[label=\alph*)]
    \item Convertir les octets 0x20 à 0x27 en caractères ASCII.
    \item Si \texttt{char *s1 = (char *)0x20}, quelle est la chaîne pointée ?
    \item Quelle est la valeur de \texttt{strlen(s1)} ?
    \item Quelle est la valeur de \texttt{sizeof("Bonjour")} ?
    \item Expliquer la différence entre \texttt{strlen} et \texttt{sizeof} pour les chaînes.
\end{enumerate}

\textbf{Partie B :} Comparer les deux déclarations :

\begin{lstlisting}[numbers=none]
char arr[] = "Hello";
char *ptr = "Hello";
\end{lstlisting}

\begin{enumerate}[label=\alph*)]
    \setcounter{enumi}{5}
    \item Quelle est la valeur de \texttt{sizeof(arr)} ?
    \item Quelle est la valeur de \texttt{sizeof(ptr)} sur une machine 64 bits ?
    \item L'instruction \texttt{arr[0] = 'X';} est-elle valide ?
    \item L'instruction \texttt{ptr[0] = 'X';} est-elle valide ? Pourquoi ?
    \item L'instruction \texttt{ptr = "World";} est-elle valide ? Et \texttt{arr = "World";} ?
\end{enumerate}

\textit{Table ASCII : 0x00=NUL, 0x41='A', 0x42='B', 0x43='C', 0x44='D', 0x55='U', 0x6A='j', 0x6E='n', 0x6F='o', 0x72='r', 0x75='u'}

\end{exercicebox}

%==============================================================================
\begin{exercicebox}{Exercice 4 : Allocation Dynamique}
%==============================================================================

\textbf{Partie A :}
\begin{lstlisting}[numbers=none]
int *a = malloc(5 * sizeof(int));
int *b = calloc(5, sizeof(int));
\end{lstlisting}

\begin{enumerate}[label=\alph*)]
    \item Combien d'octets sont alloués pour \texttt{a} ? Pour \texttt{b} ?
    \item Quelle est la différence entre les contenus initiaux de \texttt{a} et \texttt{b} ?
    \item Pourquoi est-il dangereux de lire \texttt{a[0]} juste après le \texttt{malloc} ?
    \item Écrire le code pour libérer correctement la mémoire.
\end{enumerate}

\textbf{Partie B :} Pour chaque code, identifier l'erreur et expliquer la conséquence :

\begin{lstlisting}[numbers=none]
/* Cas 1 : Use-after-free */
int *p = malloc(sizeof(int));
free(p);
*p = 42;

/* Cas 2 : Memory leak */
int *q = malloc(10 * sizeof(int));
q = malloc(20 * sizeof(int));

/* Cas 3 : Double free */
int *r = malloc(sizeof(int));
free(r);
free(r);
\end{lstlisting}

\end{exercicebox}

%==============================================================================
\begin{exercicebox}{Exercice 5 : Forensics Mémoire I (32 bits, Little-Endian)}
%==============================================================================

Soit le dump mémoire suivant :

\begin{center}
\renewcommand{\arraystretch}{1.2}
\begin{tabular}{|c|c|c|c|c|c|c|c|c|}
\hline
\textbf{Addr} & \textbf{+0} & \textbf{+1} & \textbf{+2} & \textbf{+3} & \textbf{+4} & \textbf{+5} & \textbf{+6} & \textbf{+7} \\
\hline
0x00 & 1e & 00 & 00 & 00 & 00 & 00 & 00 & 00 \\
\hline
0x08 & aa & bb & cc & dd & ee & ff & 00 & 11 \\
\hline
0x10 & 8a & 7c & 6f & 22 & 9a & 66 & 44 & 17 \\
\hline
0x18 & 33 & 77 & 6f & 6c & 66 & 62 & 79 & 74 \\
\hline
0x20 & 65 & 73 & 00 & 44 & 00 & 00 & 88 & c0 \\
\hline
0x28 & de & ad & be & ef & ca & fe & f0 & 0d \\
\hline
\end{tabular}
\end{center}

Sachant que :
\begin{lstlisting}[numbers=none]
char  *c = (char *)0x19;
short *s = (short *)0x0C;
int   *p = (int *)0x28;
\end{lstlisting}

Remplir le tableau avec le \textbf{type C} et la \textbf{valeur hexadécimale} :

\begin{center}
\renewcommand{\arraystretch}{1.5}
\begin{tabular}{|l|c|c|}
\hline
\textbf{Expression C} & \textbf{Type C} & \textbf{Valeur Hex} \\
\hline
\texttt{*s} & & \\
\hline
\texttt{*p} & & \\
\hline
\texttt{*(c+2)} & & \\
\hline
\texttt{p[1]} & & \\
\hline
\texttt{*((short *)(c-1))} & & \\
\hline
\texttt{strlen(c)} & & \\
\hline
\end{tabular}
\end{center}

\end{exercicebox}

%==============================================================================
\begin{exercicebox}{Exercice 6 : Forensics Mémoire II — Double Indirection}
%==============================================================================

Machine 32 bits, little-endian. Dump mémoire :

\begin{center}
\renewcommand{\arraystretch}{1.2}
\begin{tabular}{|c|c|c|c|c|c|c|c|c|}
\hline
\textbf{Addr} & \textbf{+0} & \textbf{+1} & \textbf{+2} & \textbf{+3} & \textbf{+4} & \textbf{+5} & \textbf{+6} & \textbf{+7} \\
\hline
0x1000 & 0A & 00 & 00 & 00 & 14 & 00 & 00 & 00 \\
\hline
0x1008 & 1E & 00 & 00 & 00 & 00 & 10 & 00 & 00 \\
\hline
0x1010 & 0C & 10 & 00 & 00 & 00 & 00 & 00 & 00 \\
\hline
\end{tabular}
\end{center}

Variables :
\begin{lstlisting}[numbers=none]
int arr[3];        /* &arr[0] = 0x1000 */
int *p;            /* &p = 0x100C */
int **pp;          /* &pp = 0x1010 */
\end{lstlisting}

\begin{enumerate}[label=\alph*)]
    \item Donner les valeurs de \texttt{arr[0]}, \texttt{arr[1]}, \texttt{arr[2]}.
    \item Quelle est la valeur de \texttt{p} ? Vers quelle adresse pointe-t-il ?
    \item Quelle est la valeur de \texttt{*p} ?
    \item Quelle est la valeur de \texttt{pp} ? Vers quelle adresse pointe-t-il ?
    \item Quelle est la valeur de \texttt{*pp} ?
    \item Quelle est la valeur de \texttt{**pp} ?
    \item Dessiner un schéma montrant les pointeurs et leurs cibles.
\end{enumerate}

\end{exercicebox}

%==============================================================================
\begin{infobox}{Rappels Essentiels}
\textbf{Endianness :}
\begin{itemize}
    \item \textbf{Little-endian} (x86) : LSB à l'adresse basse
    \item \textbf{Big-endian} : MSB à l'adresse basse
\end{itemize}

\textbf{Arithmétique des pointeurs :}
\begin{itemize}
    \item \texttt{p + n} ajoute \texttt{n × sizeof(*p)} octets
    \item \texttt{p[i]} $\equiv$ \texttt{*(p + i)}
    \item \texttt{q - p} = nombre d'éléments entre \texttt{p} et \texttt{q}
\end{itemize}

\textbf{Chaînes :}
\begin{itemize}
    \item Terminées par \texttt{'\textbackslash 0'} (0x00)
    \item \texttt{strlen} ne compte pas le terminateur
    \item \texttt{sizeof} compte le terminateur
\end{itemize}

\textbf{Valeurs suspectes en debug :}
\begin{itemize}
    \item \texttt{0xDEADBEEF}, \texttt{0xCAFEBABE} : marqueurs de debug
    \item \texttt{0xCDCDCDCD} : mémoire non initialisée (MSVC)
    \item \texttt{0xDDDDDDDD} : mémoire libérée (MSVC)
\end{itemize}
\end{infobox}

\end{document}
